% GENERATED FILE. Edit canonical lessons, then run npm run generate:books.
% canonical-source-hash: fnv1a64:2434330e5ff6ff3f
% canonical-lessons: FR-C250-etranger, FR-C250-touriste, FR-C250-adulte, FR-C250-personne-agee, FR-C250-jumeau, FR-R250-first-pass-words-from-chapters-242-245, FR-R250-second-pass-words-from-chapters-246-250

\chapter{Foreigner, Tourist, Adult, Elderly person, Twin}
\label{ch:fr-foreigner-tourist-adult-elderly-person-twin}
\hlchaptermodality{250}

\begin{chapteropening}
\textbf{By the end of this chapter:} \emph{I can say a foreigner, a tourist, an adult, an elderly person and a twin.}

Complete the last lesson of chapter 250: I can say a foreigner, a tourist, an adult, an elderly person and a twin.
\end{chapteropening}

\section[l'étranger]{l'étranger — a foreigner}

\label{lesson:FR-C250-etranger}



\begin{quote}
Before the new one: say the French for a writer, then the French for a painter.
\end{quote}

\subsection*{You'll want to know: l'étranger}
\textbf{l'étranger} — \textquotedblleft{}a foreigner\textquotedblright{}.

\begin{grammarlens}[title={What you've built}]
One more word for this chapter.
\end{grammarlens}

\subsection*{Your turn}
\begin{itemize}
\raggedright
  \item \emph{Say it:} \emph{l'étranger}
  \item \emph{Say it:} \emph{l'étranger}, once more
  \item \emph{Say it:} say \emph{l'étranger}
  \item \emph{Recall:} say the French for a writer, then the French for a painter, then say \emph{l'étranger} again
\end{itemize}

\begin{tcolorbox}[breakable,colback=teal!4,colframe=teal!35!black,title={Before you move on}]
What does l'étranger mean? (\textquotedblleft{}A foreigner\textquotedblright{}.) Say it once more.
\end{tcolorbox}
\section[le touriste]{le touriste — a tourist}

\label{lesson:FR-C250-touriste}



\begin{quote}
Before the new one: say the French for a painter, then the French for a foreigner.
\end{quote}

\subsection*{You'll want to know: le touriste}
\textbf{le touriste} — \textquotedblleft{}a tourist\textquotedblright{}.

\begin{grammarlens}[title={What you've built}]
One more word for this chapter.
\end{grammarlens}

\subsection*{Your turn}
\begin{itemize}
\raggedright
  \item \emph{Say it:} \emph{le touriste}
  \item \emph{Say it:} \emph{le touriste}, once more
  \item \emph{Say it:} say \emph{le touriste}
  \item \emph{Recall:} say the French for a painter, then the French for a foreigner, then say \emph{le touriste} again
\end{itemize}

\begin{tcolorbox}[breakable,colback=teal!4,colframe=teal!35!black,title={Before you move on}]
What does le touriste mean? (\textquotedblleft{}A tourist\textquotedblright{}.) Say it once more.
\end{tcolorbox}
\section[l'adulte]{l'adulte — an adult}

\label{lesson:FR-C250-adulte}



\begin{quote}
Before the new one: say the French for a foreigner, then the French for a tourist.
\end{quote}

\subsection*{You'll want to know: l'adulte}
\textbf{l'adulte} — \textquotedblleft{}an adult\textquotedblright{}.

\begin{grammarlens}[title={What you've built}]
One more word for this chapter.
\end{grammarlens}

\subsection*{Your turn}
\begin{itemize}
\raggedright
  \item \emph{Say it:} \emph{l'adulte}
  \item \emph{Say it:} \emph{l'adulte}, once more
  \item \emph{Say it:} say \emph{l'adulte}
  \item \emph{Recall:} say the French for a foreigner, then the French for a tourist, then say \emph{l'adulte} again
\end{itemize}

\begin{tcolorbox}[breakable,colback=teal!4,colframe=teal!35!black,title={Before you move on}]
What does l'adulte mean? (\textquotedblleft{}An adult\textquotedblright{}.) Say it once more.
\end{tcolorbox}
\section[la personne âgée]{la personne âgée — an elderly person}

\label{lesson:FR-C250-personne-agee}



\begin{quote}
Before the new one: say the French for a tourist, then the French for an adult.
\end{quote}

\subsection*{You'll want to know: la personne âgée}
\textbf{la personne âgée} — \textquotedblleft{}an elderly person\textquotedblright{}.

\begin{grammarlens}[title={What you've built}]
One more word for this chapter.
\end{grammarlens}

\subsection*{Your turn}
\begin{itemize}
\raggedright
  \item \emph{Say it:} \emph{la personne âgée}
  \item \emph{Say it:} \emph{la personne âgée}, once more
  \item \emph{Say it:} say \emph{la personne âgée}
  \item \emph{Recall:} say the French for a tourist, then the French for an adult, then say \emph{la personne âgée} again
\end{itemize}

\begin{tcolorbox}[breakable,colback=teal!4,colframe=teal!35!black,title={Before you move on}]
What does la personne âgée mean? (\textquotedblleft{}An elderly person\textquotedblright{}.) Say it once more.
\end{tcolorbox}
\section[le jumeau]{le jumeau — a twin}

\label{lesson:FR-C250-jumeau}



\begin{quote}
Before the new one: say the French for an adult, then the French for an elderly person.
\end{quote}

\subsection*{You'll want to know: le jumeau}
\textbf{le jumeau} — \textquotedblleft{}a twin\textquotedblright{}.

\begin{grammarlens}[title={What you've built}]
One more word for this chapter.
\end{grammarlens}

\subsection*{Your turn}
\begin{itemize}
\raggedright
  \item \emph{Say it:} \emph{le jumeau}
  \item \emph{Say it:} \emph{le jumeau}, once more
  \item \emph{Say it:} say \emph{le jumeau}
  \item \emph{Recall:} say the French for an adult, then the French for an elderly person, then say \emph{le jumeau} again
\end{itemize}

\begin{tcolorbox}[breakable,colback=teal!4,colframe=teal!35!black,title={Before you move on}]
What does le jumeau mean? (\textquotedblleft{}A twin\textquotedblright{}.) Say it once more.
\end{tcolorbox}
\section[(dialogue)]{Words from chapters 242-245 — a first pass}

\label{lesson:FR-R250-first-pass-words-from-chapters-242-245}



\begin{quote}
Say the words for an elderly person and a twin.
\end{quote}

\subsection*{Your turn}
Say these aloud:

\begin{itemize}
\raggedright
  \item to lose weight, to feed, to go on holiday, to joke, to tidy up
  \item a plaster, the emergency department, an emergency, a medical appointment, a diet
  \item silent, clean, hard, smooth, rough
  \item similar, fair, false, ready, free of charge
\end{itemize}

\begin{tcolorbox}[breakable,colback=teal!4,colframe=teal!35!black,title={Before you move on}]
To lose weight? (\textbf{maigrir}.) To feed? (\textbf{nourrir}.) To go on holiday? (\textbf{partir en vacances}.) To joke? (\textbf{plaisanter}.) To tidy up? (\textbf{ranger}.) A plaster? (\textbf{le pansement}.) The emergency department? (\textbf{les urgences}.) An emergency? (\textbf{l'urgence}.) A medical appointment? (\textbf{la consultation}.) A diet? (\textbf{le régime}.) Silent? (\textbf{silencieux / silencieuse}.) Clean? (\textbf{propre}.) Hard? (\textbf{dur / dure}.) Smooth? (\textbf{lisse}.) Rough? (\textbf{rugueux / rugueuse}.) Similar? (\textbf{semblable}.) Fair? (\textbf{juste}.) False? (\textbf{faux / fausse}.) Ready? (\textbf{prêt / prête}.) Free of charge? (\textbf{gratuit / gratuite}.)
\end{tcolorbox}
\section[(dialogue)]{Words from chapters 246-250 — a second pass}

\label{lesson:FR-R250-second-pass-words-from-chapters-246-250}



\begin{quote}
Say the words for generous and lazy.
\end{quote}

\subsection*{Your turn}
Say these aloud:

\begin{itemize}
\raggedright
  \item generous, lazy, hard-working, talkative, angry
  \item the bill, a tip, a stop, a motorbike, a platform
  \item a mayor, elections, a vote, politics, the economy
  \item an actor, an actress, a poem, a writer, a painter
  \item a foreigner, a tourist, an adult, an elderly person, a twin
\end{itemize}

\begin{tcolorbox}[breakable,colback=teal!4,colframe=teal!35!black,title={Before you move on}]
Generous? (\textbf{généreux / généreuse}.) Lazy? (\textbf{paresseux / paresseuse}.) Hard-working? (\textbf{travailleur / travailleuse}.) Talkative? (\textbf{bavard / bavarde}.) Angry? (\textbf{fâché / fâchée}.) The bill? (\textbf{l'addition}.) A tip? (\textbf{le pourboire}.) A stop? (\textbf{l'arrêt}.) A motorbike? (\textbf{la moto}.) A platform? (\textbf{le quai}.) A mayor? (\textbf{le maire}.) Elections? (\textbf{les élections}.) A vote? (\textbf{le vote}.) Politics? (\textbf{la politique}.) The economy? (\textbf{l'économie}.) An actor? (\textbf{l'acteur}.) An actress? (\textbf{l'actrice}.) A poem? (\textbf{le poème}.) A writer? (\textbf{l'écrivain}.) A painter? (\textbf{le peintre}.) A foreigner? (\textbf{l'étranger}.) A tourist? (\textbf{le touriste}.) An adult? (\textbf{l'adulte}.) An elderly person? (\textbf{la personne âgée}.) A twin? (\textbf{le jumeau}.)
\end{tcolorbox}
